% ---------------------------------------------------------------
%  شکل: معماری زنجیره‌ی پیشنهادی، با داده‌ی واقعی
%
%  نیازمند: tikz + positioning, arrows.meta, calc
%           (در styles/common.tex و پیش از xepersian بارگذاری شده‌اند)
%
%  استفاده:  % ---------------------------------------------------------------
%  شکل: معماری زنجیره‌ی پیشنهادی، با داده‌ی واقعی
%
%  نیازمند: tikz + positioning, arrows.meta, calc
%           (در styles/common.tex و پیش از xepersian بارگذاری شده‌اند)
%
%  استفاده:  % ---------------------------------------------------------------
%  شکل: معماری زنجیره‌ی پیشنهادی، با داده‌ی واقعی
%
%  نیازمند: tikz + positioning, arrows.meta, calc
%           (در styles/common.tex و پیش از xepersian بارگذاری شده‌اند)
%
%  استفاده:  % ---------------------------------------------------------------
%  شکل: معماری زنجیره‌ی پیشنهادی، با داده‌ی واقعی
%
%  نیازمند: tikz + positioning, arrows.meta, calc
%           (در styles/common.tex و پیش از xepersian بارگذاری شده‌اند)
%
%  استفاده:  \input{figs/pipeline_final}
%
%  ستون چپ: نوارهای تصویری figs/registration/row*.pdf که با
%  figs/registration/make_pipeline_panels.py از یک نوبت واقعی
%  VinDr-Multiphase ساخته می‌شوند. اندازه‌ی هر نوار ۱۰ × ۲٫۵۵ سانتی‌متر
%  است؛ اگر آن را در اسکریپت عوض کردی، \Rw و \Rh زیر را هم عوض کن.
%
%  ▲ قاعده‌ی مهم دوزبانه در این فایل ▲
%  متن داخل نودهای tikz جهت راست‌به‌چپ را از xepersian به ارث نمی‌برد.
%  آزمایش‌شده: \setRTL تنها، و execute at begin node، هیچ‌کدام جهت را
%  عوض نمی‌کنند؛ \rl{} عوض می‌کند ولی داخل \begin{scope} ساختار محیط را
%  می‌شکند. راه درست \Prtl زیر است (parbox با \setRTL).
%  هر نود متنی تازه‌ای که اضافه می‌کنی باید در آن بپیچد،
%  وگرنه کلمه‌به‌کلمه چپ‌به‌راست چیده می‌شود.
%  اصطلاح لاتین داخل همان parbox با \lr{...} می‌آید.
% ---------------------------------------------------------------

\begingroup
\newcommand{\Stitle}{\small\bfseries}
\newcommand{\Ssub}{\scriptsize\color{xtire}}
\newcommand{\Prtl}[2]{\parbox{#1}{\setRTL\raggedleft #2}}
\def\Rw{10cm}     % عرض نوار تصویری
\def\Rh{2.55cm}   % ارتفاع نوار تصویری
\def\Rp{2.78}     % فاصله‌ی عمودی سطرها (سانتی‌متر)

\begin{figure}[p]
\centering
\begin{tikzpicture}[
  font=\scriptsize,
  >={Stealth[length=2mm,width=1.6mm]},
  box/.style={draw=xabi, line width=.8pt, fill=white, rounded corners=2pt,
              inner xsep=2.4mm, inner ysep=1.8mm, anchor=west},
  num/.style={circle, draw=xabi, line width=.6pt, fill=white, text=xabi,
              inner sep=0pt, minimum size=4.6mm, font=\scriptsize\bfseries},
  flow/.style={->, draw=xabi, line width=.9pt},
  strip/.style={inner sep=0pt, anchor=east},
  link/.style={draw=xnokhodi, line width=.5pt},
]

% ═══════════ ستون گام‌ها (راست) ═══════════
\node[box] (s1) at (0.3, 0) {\Prtl{48.5mm}{%
  {\Stitle انتخاب سری}\\ {\Ssub برای هر نوبت و هر پویش، نازک‌ترین سری انتخاب می‌شود}}};

\node[box] (s2) at (0.3, -\Rp) {\Prtl{48.5mm}{%
  {\Stitle استانداردسازی هندسه و شدت}\\ {\Ssub جهت \lr{LPS}، بازنمونه‌برداری همسان‌گرد \عدد{۱٫۵} میلی‌متر}}};

\node[box] (s3) at (0.3, -2*\Rp) {\Prtl{48.5mm}{%
  {\Stitle بخش‌بندی اندام‌ها}\\
  {\Ssub با \lr{TotalSegmentator}؛ نقاب هم‌ترازسازی از استخوان‌ها و بافت نرم پایدار ساخته می‌شود}\\
  {\Ssub روده، دوازدهه و رکتوم کنار گذاشته می‌شوند}}};

\node[box] (s4) at (0.3, -3*\Rp) {\Prtl{48.5mm}{%
  {\Stitle برش ناحیه‌ی مشترک}\\ {\Ssub کادر قلب تا مثانه، مشترک بین دو پویش}}};

\node[box] (s5) at (0.3, -4*\Rp) {\Prtl{48.5mm}{%
  {\Stitle هم‌ترازسازی صلب کابش}\\
  {\Ssub چرخش و انتقال، حل بسته از مرکز جرم اندام‌های نقاب هم‌ترازسازی}}};

\node[box] (s6) at (0.3, -5*\Rp) {\Prtl{48.5mm}{%
  {\Stitle پالایش نلدر-مید}\\
  {\Ssub کاهش فاصله‌ی مراکز و افزایش هم‌پوشانی، بدون مشتق}\\
  {\Ssub روی \عدد{۱۲} ساختار: کبد، طحال، دو کلیه، طناب نخاعی، مهره‌های \lr{L1} تا \lr{L5} و دو عضله‌ی ایلیوپسواس}}};

\node[box] (s7) at (0.3, -6*\Rp) {\Prtl{48.5mm}{%
  {\Stitle ارزیابی}\\ {\Ssub روی \عدد{۱۱} اندام؛ تنها پانکراس و کیسه‌ی صفرا از هر دو گذر برازش کنار گذاشته شده‌اند}}};

\foreach \i/\n in {1/s1,2/s2,3/s3,4/s4,5/s5,6/s6,7/s7}
  { \node[num] at (\n.north east) {\i}; }

\foreach \a/\b in {s1/s2,s2/s3,s3/s4,s4/s5,s5/s6,s6/s7}
  { \draw[flow] (\a) -- (\b); }

% ═══════════ ستون تصویری (چپ): داده‌ی واقعی پس از هر گام ═══════════
\foreach \k/\f in {0/row1_series, 1/row2_standardize, 2/row3_segment, 3/row4_crop,
                   4/row5_pass0, 5/row6_pass1, 6/row7_eval}
  { \node[strip] at (0, -\k*\Rp)
      {\includegraphics[width=\Rw, height=\Rh]{registration/\f.pdf}}; }

\foreach \n in {s1,s2,s3,s4,s5,s6,s7}
  { \draw[link] (0.02, 0 |- \n.west) -- (\n.west); }

\end{tikzpicture}

\caption[معماری زنجیره‌ی پیشنهادی]{%
معماری زنجیره‌ی پیشنهادی. ستون راست گام‌ها و ستون چپ داده‌ی واقعی یک نوبت
\lr{\mbox{VinDr-Multiphase}} پس از همان گام است (\lr{NC} ثابت، \lr{VEN} متحرک). در
گام ۲ نوار قرمز بخشی است که در همه‌ی پویش‌ها تصویر نشده؛ در گام ۳ آبی و سبز نقاب
هم‌ترازسازی و زرد ساختارهای ویژه‌ی ارزیابی‌اند؛ در گام ۴ تصویر راست، شطرنجی دو
پویش پیش از هم‌ترازسازی است. در گام‌های ۵ و ۶ مرز اندام در \lr{NC} فیروزه‌ای و در
\lr{VEN} نارنجی است و عدد هر قاب، \lr{Dice} فرسوده پیش از هم‌ترازسازی و پس از همان
گذر است. سهم هر گام روی کل مجموعه در فصل نتایج گزارش شده است.}
\label{fig:pipeline-final}
\end{figure}
\endgroup

%
%  ستون چپ: نوارهای تصویری figs/registration/row*.pdf که با
%  figs/registration/make_pipeline_panels.py از یک نوبت واقعی
%  VinDr-Multiphase ساخته می‌شوند. اندازه‌ی هر نوار ۱۰ × ۲٫۵۵ سانتی‌متر
%  است؛ اگر آن را در اسکریپت عوض کردی، \Rw و \Rh زیر را هم عوض کن.
%
%  ▲ قاعده‌ی مهم دوزبانه در این فایل ▲
%  متن داخل نودهای tikz جهت راست‌به‌چپ را از xepersian به ارث نمی‌برد.
%  آزمایش‌شده: \setRTL تنها، و execute at begin node، هیچ‌کدام جهت را
%  عوض نمی‌کنند؛ \rl{} عوض می‌کند ولی داخل \begin{scope} ساختار محیط را
%  می‌شکند. راه درست \Prtl زیر است (parbox با \setRTL).
%  هر نود متنی تازه‌ای که اضافه می‌کنی باید در آن بپیچد،
%  وگرنه کلمه‌به‌کلمه چپ‌به‌راست چیده می‌شود.
%  اصطلاح لاتین داخل همان parbox با \lr{...} می‌آید.
% ---------------------------------------------------------------

\begingroup
\newcommand{\Stitle}{\small\bfseries}
\newcommand{\Ssub}{\scriptsize\color{xtire}}
\newcommand{\Prtl}[2]{\parbox{#1}{\setRTL\raggedleft #2}}
\def\Rw{10cm}     % عرض نوار تصویری
\def\Rh{2.55cm}   % ارتفاع نوار تصویری
\def\Rp{2.78}     % فاصله‌ی عمودی سطرها (سانتی‌متر)

\begin{figure}[p]
\centering
\begin{tikzpicture}[
  font=\scriptsize,
  >={Stealth[length=2mm,width=1.6mm]},
  box/.style={draw=xabi, line width=.8pt, fill=white, rounded corners=2pt,
              inner xsep=2.4mm, inner ysep=1.8mm, anchor=west},
  num/.style={circle, draw=xabi, line width=.6pt, fill=white, text=xabi,
              inner sep=0pt, minimum size=4.6mm, font=\scriptsize\bfseries},
  flow/.style={->, draw=xabi, line width=.9pt},
  strip/.style={inner sep=0pt, anchor=east},
  link/.style={draw=xnokhodi, line width=.5pt},
]

% ═══════════ ستون گام‌ها (راست) ═══════════
\node[box] (s1) at (0.3, 0) {\Prtl{48.5mm}{%
  {\Stitle انتخاب سری}\\ {\Ssub برای هر نوبت و هر پویش، نازک‌ترین سری انتخاب می‌شود}}};

\node[box] (s2) at (0.3, -\Rp) {\Prtl{48.5mm}{%
  {\Stitle استانداردسازی هندسه و شدت}\\ {\Ssub جهت \lr{LPS}، بازنمونه‌برداری همسان‌گرد \عدد{۱٫۵} میلی‌متر}}};

\node[box] (s3) at (0.3, -2*\Rp) {\Prtl{48.5mm}{%
  {\Stitle بخش‌بندی اندام‌ها}\\
  {\Ssub با \lr{TotalSegmentator}؛ نقاب هم‌ترازسازی از استخوان‌ها و بافت نرم پایدار ساخته می‌شود}\\
  {\Ssub روده، دوازدهه و رکتوم کنار گذاشته می‌شوند}}};

\node[box] (s4) at (0.3, -3*\Rp) {\Prtl{48.5mm}{%
  {\Stitle برش ناحیه‌ی مشترک}\\ {\Ssub کادر قلب تا مثانه، مشترک بین دو پویش}}};

\node[box] (s5) at (0.3, -4*\Rp) {\Prtl{48.5mm}{%
  {\Stitle هم‌ترازسازی صلب کابش}\\
  {\Ssub چرخش و انتقال، حل بسته از مرکز جرم اندام‌های نقاب هم‌ترازسازی}}};

\node[box] (s6) at (0.3, -5*\Rp) {\Prtl{48.5mm}{%
  {\Stitle پالایش نلدر-مید}\\
  {\Ssub کاهش فاصله‌ی مراکز و افزایش هم‌پوشانی، بدون مشتق}\\
  {\Ssub روی \عدد{۱۲} ساختار: کبد، طحال، دو کلیه، طناب نخاعی، مهره‌های \lr{L1} تا \lr{L5} و دو عضله‌ی ایلیوپسواس}}};

\node[box] (s7) at (0.3, -6*\Rp) {\Prtl{48.5mm}{%
  {\Stitle ارزیابی}\\ {\Ssub روی \عدد{۱۱} اندام؛ تنها پانکراس و کیسه‌ی صفرا از هر دو گذر برازش کنار گذاشته شده‌اند}}};

\foreach \i/\n in {1/s1,2/s2,3/s3,4/s4,5/s5,6/s6,7/s7}
  { \node[num] at (\n.north east) {\i}; }

\foreach \a/\b in {s1/s2,s2/s3,s3/s4,s4/s5,s5/s6,s6/s7}
  { \draw[flow] (\a) -- (\b); }

% ═══════════ ستون تصویری (چپ): داده‌ی واقعی پس از هر گام ═══════════
\foreach \k/\f in {0/row1_series, 1/row2_standardize, 2/row3_segment, 3/row4_crop,
                   4/row5_pass0, 5/row6_pass1, 6/row7_eval}
  { \node[strip] at (0, -\k*\Rp)
      {\includegraphics[width=\Rw, height=\Rh]{registration/\f.pdf}}; }

\foreach \n in {s1,s2,s3,s4,s5,s6,s7}
  { \draw[link] (0.02, 0 |- \n.west) -- (\n.west); }

\end{tikzpicture}

\caption[معماری زنجیره‌ی پیشنهادی]{%
معماری زنجیره‌ی پیشنهادی. ستون راست گام‌ها و ستون چپ داده‌ی واقعی یک نوبت
\lr{\mbox{VinDr-Multiphase}} پس از همان گام است (\lr{NC} ثابت، \lr{VEN} متحرک). در
گام ۲ نوار قرمز بخشی است که در همه‌ی پویش‌ها تصویر نشده؛ در گام ۳ آبی و سبز نقاب
هم‌ترازسازی و زرد ساختارهای ویژه‌ی ارزیابی‌اند؛ در گام ۴ تصویر راست، شطرنجی دو
پویش پیش از هم‌ترازسازی است. در گام‌های ۵ و ۶ مرز اندام در \lr{NC} فیروزه‌ای و در
\lr{VEN} نارنجی است و عدد هر قاب، \lr{Dice} فرسوده پیش از هم‌ترازسازی و پس از همان
گذر است. سهم هر گام روی کل مجموعه در فصل نتایج گزارش شده است.}
\label{fig:pipeline-final}
\end{figure}
\endgroup

%
%  ستون چپ: نوارهای تصویری figs/registration/row*.pdf که با
%  figs/registration/make_pipeline_panels.py از یک نوبت واقعی
%  VinDr-Multiphase ساخته می‌شوند. اندازه‌ی هر نوار ۱۰ × ۲٫۵۵ سانتی‌متر
%  است؛ اگر آن را در اسکریپت عوض کردی، \Rw و \Rh زیر را هم عوض کن.
%
%  ▲ قاعده‌ی مهم دوزبانه در این فایل ▲
%  متن داخل نودهای tikz جهت راست‌به‌چپ را از xepersian به ارث نمی‌برد.
%  آزمایش‌شده: \setRTL تنها، و execute at begin node، هیچ‌کدام جهت را
%  عوض نمی‌کنند؛ \rl{} عوض می‌کند ولی داخل \begin{scope} ساختار محیط را
%  می‌شکند. راه درست \Prtl زیر است (parbox با \setRTL).
%  هر نود متنی تازه‌ای که اضافه می‌کنی باید در آن بپیچد،
%  وگرنه کلمه‌به‌کلمه چپ‌به‌راست چیده می‌شود.
%  اصطلاح لاتین داخل همان parbox با \lr{...} می‌آید.
% ---------------------------------------------------------------

\begingroup
\newcommand{\Stitle}{\small\bfseries}
\newcommand{\Ssub}{\scriptsize\color{xtire}}
\newcommand{\Prtl}[2]{\parbox{#1}{\setRTL\raggedleft #2}}
\def\Rw{10cm}     % عرض نوار تصویری
\def\Rh{2.55cm}   % ارتفاع نوار تصویری
\def\Rp{2.78}     % فاصله‌ی عمودی سطرها (سانتی‌متر)

\begin{figure}[p]
\centering
\begin{tikzpicture}[
  font=\scriptsize,
  >={Stealth[length=2mm,width=1.6mm]},
  box/.style={draw=xabi, line width=.8pt, fill=white, rounded corners=2pt,
              inner xsep=2.4mm, inner ysep=1.8mm, anchor=west},
  num/.style={circle, draw=xabi, line width=.6pt, fill=white, text=xabi,
              inner sep=0pt, minimum size=4.6mm, font=\scriptsize\bfseries},
  flow/.style={->, draw=xabi, line width=.9pt},
  strip/.style={inner sep=0pt, anchor=east},
  link/.style={draw=xnokhodi, line width=.5pt},
]

% ═══════════ ستون گام‌ها (راست) ═══════════
\node[box] (s1) at (0.3, 0) {\Prtl{48.5mm}{%
  {\Stitle انتخاب سری}\\ {\Ssub برای هر نوبت و هر پویش، نازک‌ترین سری انتخاب می‌شود}}};

\node[box] (s2) at (0.3, -\Rp) {\Prtl{48.5mm}{%
  {\Stitle استانداردسازی هندسه و شدت}\\ {\Ssub جهت \lr{LPS}، بازنمونه‌برداری همسان‌گرد \عدد{۱٫۵} میلی‌متر}}};

\node[box] (s3) at (0.3, -2*\Rp) {\Prtl{48.5mm}{%
  {\Stitle بخش‌بندی اندام‌ها}\\
  {\Ssub با \lr{TotalSegmentator}؛ نقاب هم‌ترازسازی از استخوان‌ها و بافت نرم پایدار ساخته می‌شود}\\
  {\Ssub روده، دوازدهه و رکتوم کنار گذاشته می‌شوند}}};

\node[box] (s4) at (0.3, -3*\Rp) {\Prtl{48.5mm}{%
  {\Stitle برش ناحیه‌ی مشترک}\\ {\Ssub کادر قلب تا مثانه، مشترک بین دو پویش}}};

\node[box] (s5) at (0.3, -4*\Rp) {\Prtl{48.5mm}{%
  {\Stitle هم‌ترازسازی صلب کابش}\\
  {\Ssub چرخش و انتقال، حل بسته از مرکز جرم اندام‌های نقاب هم‌ترازسازی}}};

\node[box] (s6) at (0.3, -5*\Rp) {\Prtl{48.5mm}{%
  {\Stitle پالایش نلدر-مید}\\
  {\Ssub کاهش فاصله‌ی مراکز و افزایش هم‌پوشانی، بدون مشتق}\\
  {\Ssub روی \عدد{۱۲} ساختار: کبد، طحال، دو کلیه، طناب نخاعی، مهره‌های \lr{L1} تا \lr{L5} و دو عضله‌ی ایلیوپسواس}}};

\node[box] (s7) at (0.3, -6*\Rp) {\Prtl{48.5mm}{%
  {\Stitle ارزیابی}\\ {\Ssub روی \عدد{۱۱} اندام؛ تنها پانکراس و کیسه‌ی صفرا از هر دو گذر برازش کنار گذاشته شده‌اند}}};

\foreach \i/\n in {1/s1,2/s2,3/s3,4/s4,5/s5,6/s6,7/s7}
  { \node[num] at (\n.north east) {\i}; }

\foreach \a/\b in {s1/s2,s2/s3,s3/s4,s4/s5,s5/s6,s6/s7}
  { \draw[flow] (\a) -- (\b); }

% ═══════════ ستون تصویری (چپ): داده‌ی واقعی پس از هر گام ═══════════
\foreach \k/\f in {0/row1_series, 1/row2_standardize, 2/row3_segment, 3/row4_crop,
                   4/row5_pass0, 5/row6_pass1, 6/row7_eval}
  { \node[strip] at (0, -\k*\Rp)
      {\includegraphics[width=\Rw, height=\Rh]{registration/\f.pdf}}; }

\foreach \n in {s1,s2,s3,s4,s5,s6,s7}
  { \draw[link] (0.02, 0 |- \n.west) -- (\n.west); }

\end{tikzpicture}

\caption[معماری زنجیره‌ی پیشنهادی]{%
معماری زنجیره‌ی پیشنهادی. ستون راست گام‌ها و ستون چپ داده‌ی واقعی یک نوبت
\lr{\mbox{VinDr-Multiphase}} پس از همان گام است (\lr{NC} ثابت، \lr{VEN} متحرک). در
گام ۲ نوار قرمز بخشی است که در همه‌ی پویش‌ها تصویر نشده؛ در گام ۳ آبی و سبز نقاب
هم‌ترازسازی و زرد ساختارهای ویژه‌ی ارزیابی‌اند؛ در گام ۴ تصویر راست، شطرنجی دو
پویش پیش از هم‌ترازسازی است. در گام‌های ۵ و ۶ مرز اندام در \lr{NC} فیروزه‌ای و در
\lr{VEN} نارنجی است و عدد هر قاب، \lr{Dice} فرسوده پیش از هم‌ترازسازی و پس از همان
گذر است. سهم هر گام روی کل مجموعه در فصل نتایج گزارش شده است.}
\label{fig:pipeline-final}
\end{figure}
\endgroup

%
%  ستون چپ: نوارهای تصویری figs/registration/row*.pdf که با
%  figs/registration/make_pipeline_panels.py از یک نوبت واقعی
%  VinDr-Multiphase ساخته می‌شوند. اندازه‌ی هر نوار ۱۰ × ۲٫۵۵ سانتی‌متر
%  است؛ اگر آن را در اسکریپت عوض کردی، \Rw و \Rh زیر را هم عوض کن.
%
%  ▲ قاعده‌ی مهم دوزبانه در این فایل ▲
%  متن داخل نودهای tikz جهت راست‌به‌چپ را از xepersian به ارث نمی‌برد.
%  آزمایش‌شده: \setRTL تنها، و execute at begin node، هیچ‌کدام جهت را
%  عوض نمی‌کنند؛ \rl{} عوض می‌کند ولی داخل \begin{scope} ساختار محیط را
%  می‌شکند. راه درست \Prtl زیر است (parbox با \setRTL).
%  هر نود متنی تازه‌ای که اضافه می‌کنی باید در آن بپیچد،
%  وگرنه کلمه‌به‌کلمه چپ‌به‌راست چیده می‌شود.
%  اصطلاح لاتین داخل همان parbox با \lr{...} می‌آید.
% ---------------------------------------------------------------

\begingroup
\newcommand{\Stitle}{\small\bfseries}
\newcommand{\Ssub}{\scriptsize\color{xtire}}
\newcommand{\Prtl}[2]{\parbox{#1}{\setRTL\raggedleft #2}}
\def\Rw{10cm}     % عرض نوار تصویری
\def\Rh{2.55cm}   % ارتفاع نوار تصویری
\def\Rp{2.78}     % فاصله‌ی عمودی سطرها (سانتی‌متر)

\begin{figure}[p]
\centering
\begin{tikzpicture}[
  font=\scriptsize,
  >={Stealth[length=2mm,width=1.6mm]},
  box/.style={draw=xabi, line width=.8pt, fill=white, rounded corners=2pt,
              inner xsep=2.4mm, inner ysep=1.8mm, anchor=west},
  num/.style={circle, draw=xabi, line width=.6pt, fill=white, text=xabi,
              inner sep=0pt, minimum size=4.6mm, font=\scriptsize\bfseries},
  flow/.style={->, draw=xabi, line width=.9pt},
  strip/.style={inner sep=0pt, anchor=east},
  link/.style={draw=xnokhodi, line width=.5pt},
]

% ═══════════ ستون گام‌ها (راست) ═══════════
\node[box] (s1) at (0.3, 0) {\Prtl{48.5mm}{%
  {\Stitle انتخاب سری}\\ {\Ssub برای هر نوبت و هر پویش، نازک‌ترین سری انتخاب می‌شود}}};

\node[box] (s2) at (0.3, -\Rp) {\Prtl{48.5mm}{%
  {\Stitle استانداردسازی هندسه و شدت}\\ {\Ssub جهت \lr{LPS}، بازنمونه‌برداری همسان‌گرد \عدد{۱٫۵} میلی‌متر}}};

\node[box] (s3) at (0.3, -2*\Rp) {\Prtl{48.5mm}{%
  {\Stitle بخش‌بندی اندام‌ها}\\
  {\Ssub با \lr{TotalSegmentator}؛ نقاب هم‌ترازسازی از استخوان‌ها و بافت نرم پایدار ساخته می‌شود}\\
  {\Ssub روده، دوازدهه و رکتوم کنار گذاشته می‌شوند}}};

\node[box] (s4) at (0.3, -3*\Rp) {\Prtl{48.5mm}{%
  {\Stitle برش ناحیه‌ی مشترک}\\ {\Ssub کادر قلب تا مثانه، مشترک بین دو پویش}}};

\node[box] (s5) at (0.3, -4*\Rp) {\Prtl{48.5mm}{%
  {\Stitle هم‌ترازسازی صلب کابش}\\
  {\Ssub چرخش و انتقال، حل بسته از مرکز جرم اندام‌های نقاب هم‌ترازسازی}}};

\node[box] (s6) at (0.3, -5*\Rp) {\Prtl{48.5mm}{%
  {\Stitle پالایش نلدر-مید}\\
  {\Ssub کاهش فاصله‌ی مراکز و افزایش هم‌پوشانی، بدون مشتق}\\
  {\Ssub روی \عدد{۱۲} ساختار: کبد، طحال، دو کلیه، طناب نخاعی، مهره‌های \lr{L1} تا \lr{L5} و دو عضله‌ی ایلیوپسواس}}};

\node[box] (s7) at (0.3, -6*\Rp) {\Prtl{48.5mm}{%
  {\Stitle ارزیابی}\\ {\Ssub روی \عدد{۱۱} اندام؛ تنها پانکراس و کیسه‌ی صفرا از هر دو گذر برازش کنار گذاشته شده‌اند}}};

\foreach \i/\n in {1/s1,2/s2,3/s3,4/s4,5/s5,6/s6,7/s7}
  { \node[num] at (\n.north east) {\i}; }

\foreach \a/\b in {s1/s2,s2/s3,s3/s4,s4/s5,s5/s6,s6/s7}
  { \draw[flow] (\a) -- (\b); }

% ═══════════ ستون تصویری (چپ): داده‌ی واقعی پس از هر گام ═══════════
\foreach \k/\f in {0/row1_series, 1/row2_standardize, 2/row3_segment, 3/row4_crop,
                   4/row5_pass0, 5/row6_pass1, 6/row7_eval}
  { \node[strip] at (0, -\k*\Rp)
      {\includegraphics[width=\Rw, height=\Rh]{registration/\f.pdf}}; }

\foreach \n in {s1,s2,s3,s4,s5,s6,s7}
  { \draw[link] (0.02, 0 |- \n.west) -- (\n.west); }

\end{tikzpicture}

\caption[معماری زنجیره‌ی پیشنهادی]{%
معماری زنجیره‌ی پیشنهادی. ستون راست گام‌ها و ستون چپ داده‌ی واقعی یک نوبت
\lr{\mbox{VinDr-Multiphase}} پس از همان گام است (\lr{NC} ثابت، \lr{VEN} متحرک). در
گام ۲ نوار قرمز بخشی است که در همه‌ی پویش‌ها تصویر نشده؛ در گام ۳ آبی و سبز نقاب
هم‌ترازسازی و زرد ساختارهای ویژه‌ی ارزیابی‌اند؛ در گام ۴ تصویر راست، شطرنجی دو
پویش پیش از هم‌ترازسازی است. در گام‌های ۵ و ۶ مرز اندام در \lr{NC} فیروزه‌ای و در
\lr{VEN} نارنجی است و عدد هر قاب، \lr{Dice} فرسوده پیش از هم‌ترازسازی و پس از همان
گذر است. سهم هر گام روی کل مجموعه در فصل نتایج گزارش شده است.}
\label{fig:pipeline-final}
\end{figure}
\endgroup
